\documentclass[10pt,a4paper]{amsart}
\usepackage{amsmath,amssymb,amsthm,graphicx}
\usepackage{verbatim}
\newenvironment{reference}{\comment}{\endcomment}
%\newenvironment{reference}{}{}
\newenvironment{slogan}{\comment}{\endcomment}
\newenvironment{history}{\comment}{\endcomment}

% For commutative diagrams we use Xy-pic
\usepackage[all]{xy}

% We use 2cell for 2-commutative diagrams.
\xyoption{2cell}
% Theorem environments.
%
\theoremstyle{plain}
\newtheorem{theorem}[subsection]{Theorem}
\newtheorem{proposition}[subsection]{Proposition}
\newtheorem{lemma}[subsection]{Lemma}

\theoremstyle{definition}
\newtheorem{definition}[subsection]{Definition}
\newtheorem{example}[subsection]{Example}
\newtheorem{exercise}[subsection]{Exercise}
\newtheorem{situation}[subsection]{Situation}

\theoremstyle{remark}
\newtheorem{remark}[subsection]{Remark}
\newtheorem{remarks}[subsection]{Remarks}

\numberwithin{equation}{subsection}

% Macros
%
\def\lim{\mathop{\mathrm{lim}}\nolimits}
\def\colim{\mathop{\mathrm{colim}}\nolimits}
\def\Spec{\mathop{\mathrm{Spec}}}
\def\Hom{\mathop{\mathrm{Hom}}\nolimits}
\def\Ext{\mathop{\mathrm{Ext}}\nolimits}
\def\SheafHom{\mathop{\mathcal{H}\!\mathit{om}}\nolimits}
\def\SheafExt{\mathop{\mathcal{E}\!\mathit{xt}}\nolimits}
\def\Sch{\mathit{Sch}}
\def\Mor{\mathop{\mathrm{Mor}}\nolimits}
\def\Ob{\mathop{\mathrm{Ob}}\nolimits}
\def\Sh{\mathop{\mathit{Sh}}\nolimits}
\def\NL{\mathop{N\!L}\nolimits}
\def\CH{\mathop{\mathrm{CH}}\nolimits}
\def\proetale{{pro\text{-}\acute{e}tale}}
\def\etale{{\acute{e}tale}}
\def\QCoh{\mathit{QCoh}}
\def\Ker{\mathop{\mathrm{Ker}}}
\def\Im{\mathop{\mathrm{Im}}}
\def\Coker{\mathop{\mathrm{Coker}}}
\def\Coim{\mathop{\mathrm{Coim}}}

% Boxtimes
%
\DeclareMathSymbol{\boxtimes}{\mathbin}{AMSa}{"02}

%
% Macros for moduli stacks/spaces
%
\def\QCohstack{\mathcal{QC}\!\mathit{oh}}
\def\Cohstack{\mathcal{C}\!\mathit{oh}}
\def\Spacesstack{\mathcal{S}\!\mathit{paces}}
\def\Quotfunctor{\mathrm{Quot}}
\def\Hilbfunctor{\mathrm{Hilb}}
\def\Curvesstack{\mathcal{C}\!\mathit{urves}}
\def\Polarizedstack{\mathcal{P}\!\mathit{olarized}}
\def\Complexesstack{\mathcal{C}\!\mathit{omplexes}}
% \Pic is the operator that assigns to X its picard group, usage \Pic(X)
% \Picardstack_{X/B} denotes the Picard stack of X over B
% \Picardfunctor_{X/B} denotes the Picard functor of X over B
\def\Pic{\mathop{\mathrm{Pic}}\nolimits}
\def\Picardstack{\mathcal{P}\!\mathit{ic}}
\def\Picardfunctor{\mathrm{Pic}}
\def\Deformationcategory{\mathcal{D}\!\mathit{ef}}
\usepackage{cleveref}
\crefname{originalreference}{original reference}{original references}
\Crefname{originalreference}{Original reference}{Original references}
\makeatletter
\newlabel{lemma-base}{{0FHE}{source}{}{sourcepointer-1}{}}
\newlabel{lemma-base@cref}{{[originalreference][0][]0FHE}{[1][0][]source}{}{}{}}
\makeatother
\crefname{originallemma}{original lemma}{original lemmas}
\Crefname{originallemma}{Original lemma}{Original lemmas}
\makeatletter
\newlabel{varieties-lemma-smooth-separable-closed-points-dense}{{056U}{source}{}{sourcepointer-2}{}}
\newlabel{varieties-lemma-smooth-separable-closed-points-dense@cref}{{[originallemma][0][]056U}{[1][0][]source}{}{}{}}
\makeatother
\title{The fundamental class of a nonempty smooth projective scheme and vanishing for the empty scheme}
\author{The Stacks Project Authors}
\date{stacks-project:weil.tex:local-snapshot-20261012; GNU FDL 1.2 or later}
\begin{document}
\maketitle
\noindent\textit{Editorial note.} All delivered mathematical text is the author's own, copied from the pinned Stacks Project chapter file weil.tex; no statement or proof step was rewritten, reordered or removed.\par\medskip
\noindent\textit{Editorial note.} The counted proof is the author's complete argument: the identification of the fundamental class as the pullback of the class of the base, the nonzero class obtained from a closed point with separable residue field, and the product argument for the empty scheme.\par\medskip
\noindent\textit{Editorial note.} The label lemma-unit names the counted statement, and the proof environment immediately following it in the source is delivered verbatim as its proof.\par\medskip
\noindent\textit{Editorial note.} The counted statement is carried by the author's own lemma environment in the chapter source.\par\medskip
\noindent\textit{Editorial note.} Pre-open gate: python operations/agent-authoring/tools/precard\_gate.py --source weil.tex --label lemma-unit --start 2799 --end 2845 --local-id 28129 returned RESULT=OK (L28129 free, occupancy FREE, public interval clear). Pre-assembly tag self-check: canonical tag resolved and not present in reference/public-items.jsonl.\par\medskip
\noindent\textit{Editorial note.} The statement and proof are copied from the pinned chapter file; the cited inputs are kept as external pointers with their canonical Stacks tags.\par\medskip
\noindent\textit{Allowed standard prerequisites (editorial).}\begin{itemize}
\item The cohomology of the base field and the cycle class of its spectrum. Same chapter, the computation of the cohomology of the spectrum of the base field. Identifies the cycle class of the spectrum of the base field with one, which is the starting point for the computation of the fundamental class.
\item The cycle class map is compatible with pullback. Axiom (C) on the compatibility of the cycle class map with pullbacks and with proper pushforwards. Shows that the fundamental class of a scheme is the pullback of the class of the base and that pushforwards of classes behave as expected.
\item Existence of closed points with separable residue field and the shape of the cohomology of the empty scheme. Varieties, the lemma on the density of closed points with separable residue field for a smooth scheme, and axiom (B) on the behaviour under disjoint unions and base change. Produces a closed point computing a nonzero class, and determines the cohomology of the empty scheme by an argument with the product structure.
\end{itemize}\medskip
\ifdefined\hypertarget\hypertarget{sourcepointer-1}{}\fi
\noindent\textit{Source reference.} Original reference 0FHE: sources/stacks-weil/pinned-original-source.tex, source lines 2771--2771. Prerequisite pointer invoked at proof line(s) 2811; it names a standard fact of the same chapter that the counted argument uses, and every step of the counted proof built on it is retained verbatim in the delivered body.\par\medskip
\ifdefined\hypertarget\hypertarget{sourcepointer-2}{}\fi
\noindent\textit{Source reference.} Original lemma 056U: sources/stacks-chow/tags.txt, source lines 6722--6722. Prerequisite pointer invoked at proof line(s) 2820; it names a standard fact of a sibling chapter that the counted argument uses, and every step of the counted proof built on it is retained verbatim in the delivered body.\par\medskip
\begin{lemma}
\label{lemma-unit}
Assume given (D0), (D1), (D2), and (D3) satisfying (A), (B), and (C).
Let $X$ be a smooth projective scheme over $k$.
If $X = \emptyset$, then $H^*(X) = 0$.
If $X$ is nonempty, then $\gamma([X]) = 1$ and $1 \not = 0$ in $H^0(X)$.
\end{lemma}

\begin{proof}
First assume $X$ is nonempty.
Observe that $[X]$ is the pullback of $[\Spec(k)]$ by the structure morphism
$p : X \to \Spec(k)$. Hence we get $\gamma([X]) = 1$ by axiom (C)(a)
and Lemma \ref{lemma-base}. Let $X' \subset X$ be an irreducible component.
By functoriality it suffices to show $1 \not = 0$ in $H^0(X')$.
Thus we may and do assume $X$ is irreducible, and in particular
nonempty and equidimensional, say of dimension $d$.
To see that $1 \not = 0$ it suffices to show that $H^*(X)$ is nonzero.

\medskip\noindent
Let $x \in X$ be a closed point whose residue field $k'$
is separable over $k$, see
Varieties, Lemma \ref{varieties-lemma-smooth-separable-closed-points-dense}.
Let $i : \Spec(k') \to X$ be the inclusion morphism.
Denote  $p : X \to \Spec(k)$ is the structure morphism.
Observe that
$p_*i_*[\Spec(k')] = [k' : k][\Spec(k)]$ in $\CH_0(\Spec(k))$.
Using axiom (C)(b) twice and Lemma \ref{lemma-base}
we conclude that
$$
p_*i_*\gamma([\Spec(k')]) = \gamma([k' : k][\Spec(k)]) = [k' : k]
\in F = H^0(\Spec(k))
$$
is nonzero. Thus $i_*\gamma([\Spec(k)]) \in H^{2d}(X)(d)$ is nonzero
(because it maps to something nonzero via $p_*$). This concludes the proof
in case $X$ is nonempty.

\medskip\noindent
Finally, we consider the case of the empty scheme. Axiom (B)(a) gives
$H^*(\emptyset) \otimes H^*(\emptyset) = H^*(\emptyset)$ and
we get that $H^*(\emptyset)$ is either zero or $1$-dimensional
in degree $0$. Then axiom (B)(a) again shows that
$H^*(\emptyset) \otimes H^*(X) = H^*(\emptyset)$ for
all smooth projective schemes $X$ over $k$. Using axiom (A)(b)
and the nonvanishing of $H^0(X)$ we've seen above
we find that $H^*(X)$ is nonzero in at least two degrees
if $\dim(X) > 0$. This then forces $H^*(\emptyset)$ to be zero.
\end{proof}

\end{document}
